\section{Introduction}

Language models are writing more and more of the code that gets deployed, and
often nobody reads that code line by line before it ships. If a model were
misaligned, this would give it an easy way to cause harm. It could write a
solution that passes all of the developer's tests but gives a wrong answer on
one rare input that it picked, i.e. a backdoor. AI control studies how models
like this could still be used safely without assuming they can be trusted
\citep{greenblatt2024aicontrol}. In the standard setup an untrusted model writes the code
and a weaker trusted model gives each solution a suspicion score. The highest
scoring solutions are sent to human auditors, although in practice only a small
fraction, 2\% in \citet{greenblatt2024aicontrol}, can be audited.

A standard baseline in this setting is trusted monitoring
\citep{greenblatt2024aicontrol, bhatt2025ctrlz, terekhov2025controltax}, in which a weaker trusted
model reads each submitted solution and assigns a suspicion score. Safety then
depends on the monitor identifying sabotage under a limited audit budget. In the
concentrated BigCodeBench backdooring setting, \citet{schaeffer2026} show that
eliciting selective attacks reduces safety from 99\% to 59\% at a 0.5\% audit
budget. Testing could be a useful alternative
here. To catch a backdoor, a test does not have to understand the code, it just
has to find an input where the output is wrong. Writing extra tests has come up
before as a possible blue team strategy, and some control protocols already
generate exploit inputs or consistency checks to flag backdoors, but we are not
aware of work that evaluates trusted-model test generation on its own as a
backdoor detector and benchmarks it against a code-reading monitor.

We look at whether tests written by a trusted model can catch backdoors that a
code-reading monitor misses. Our reasoning was that checking an answer is often
easier than coming up with it. For a competitive programming problem, a model
might struggle to write a correct solution but still be able to list properties
that any correct solution has to satisfy, and these can be checked on lots of
generated inputs \citep{claessen2000}. If successful, these tests could provide additional detection value beyond the
monitor's suspicion scores. We expected this to be especially true when the test writer is only
given the problem statement and never sees the code.

We evaluate trusted-model test generation on all 1{,}388 eligible APPS problems
\citep{hendrycks2021apps}, using backdoored solutions from the Control Tax work
\citep{terekhov2025controltax}. The trusted model is DeepSeek-V3.2 at low
reasoning effort. We compare property tests, which may relate outputs across
inputs, with plain tests, which check a single output, and a protocol in which
the same model at high reasoning effort reads the submission and plans tests that
the trusted model implements without seeing the code. We evaluate each as a
standalone detector and as evidence given to a code-reading monitor, measure
safety with the auditing model of \citet{greenblatt2024aicontrol}, and run every
protocol with defer to trusted, the rule of that work's baseline.

To understand why the tests miss, we measure whether their inputs reach each
backdoor at all, and in a smaller study on 50 hard APPS problems we evaluate two
changes to the test writer (few-shot examples of good and bad tests, and a critic
that drops weak tests).

With defer to trusted, giving the monitor the trusted model's own tests does not
raise safety at equal usefulness: at a usefulness of 0.95, trusted monitoring and
monitoring with tests all reach about 0.73. Only tests planned by the untrusted
model raise it, to 0.82, without statistical certainty. The tests do sharpen the
top of the monitor's ranking, detecting significantly more backdoors at strict
false-positive rates, which raises safety in point estimate when nothing is
deferred, but not once the monitor can defer. Tests alone are a weak detector, and
not for lack of reach: the trusted model's inputs reach 87\% of backdoors and a
test fails on nearly every backdoor reached, but tests also fail on 28--49\% of
honest code.

Our contributions are:
\begin{itemize}
  \item We evaluate trusted-model test generation as a backdoor detector against
  a code-reading monitor on every eligible APPS problem, alone and as evidence
  for the monitor, under the standard auditing model with defer to trusted and
  with paired bootstrap intervals.
  \item We show that with defer to trusted, the trusted model's tests do not
  raise the monitor's safety at equal usefulness, although they improve its
  ranking at strict thresholds.
  \item We show that the tests reach most backdoors but cannot separate them from
  correct code: their failures come mostly from assertions the specification does
  not entail. We call this bottleneck the specification-entailment wall.
\end{itemize}
